\documentclass[11pt]{article}
% ===========================================================================
%  Shared preamble for all MPM2D worksheets — input right after \documentclass.
%  (Files beginning with "_" are skipped by mpm2d-worksheet-publish.mjs.)
% ===========================================================================
\usepackage[letterpaper,top=2.2cm,bottom=1.9cm,left=1.6cm,right=1.6cm,headheight=28pt,footskip=24pt]{geometry}
\usepackage{amsmath,amssymb}
\usepackage[T1]{fontenc}
\usepackage[utf8]{inputenc}
\usepackage{lmodern}
\usepackage{xcolor}
\usepackage[most]{tcolorbox}
\usepackage{enumitem}
\usepackage{multicol}
\usepackage{fancyhdr}
\usepackage{tikz}
\usetikzlibrary{calc}
\usepackage{pgfplots}
\pgfplotsset{compat=1.18}

\definecolor{exblue}{HTML}{4A90E2}
\definecolor{exbg}{HTML}{E6F3FF}
\definecolor{qorange}{HTML}{E69138}
\definecolor{qbg}{HTML}{FFF7CC}
\definecolor{keygray}{HTML}{475569}
\definecolor{keybg}{HTML}{F1F5F9}
\definecolor{iagreen}{HTML}{14653B}
\definecolor{titleblue}{HTML}{1D4ED8}
% Rotating palette so each worked example gets its own colour.
\definecolor{ex0}{HTML}{2563A0}\definecolor{ex1}{HTML}{3B7D3B}\definecolor{ex2}{HTML}{8A5A00}
\definecolor{ex3}{HTML}{A3327A}\definecolor{ex4}{HTML}{6D28A3}\definecolor{ex5}{HTML}{B08900}
\definecolor{ex6}{HTML}{0E7490}\definecolor{ex7}{HTML}{9A3412}\definecolor{ex8}{HTML}{15803D}

\pagestyle{fancy}
\fancyhf{}
\fancyhead[L]{\footnotesize NAME: \underline{\hspace{3.4cm}}}
\fancyhead[C]{\textbf{Integration Academy}\\[-2pt]{\scriptsize Math Simplified \textperiodcentered{} Success Amplified}}
\fancyhead[R]{\footnotesize MPM2D}
\fancyfoot[L]{\scriptsize dr.merie@IntegrationAcademy.ca}
\fancyfoot[C]{\scriptsize\thepage}
\fancyfoot[R]{\scriptsize IntegrationAcademy.ca}
\renewcommand{\headrulewidth}{1.2pt}
\renewcommand{\footrulewidth}{0.4pt}
\setlength{\parindent}{0pt}
\setlength{\parskip}{4pt}

% Examples are NOT breakable, so each example + its solution stays on one page.
\newtcolorbox{example}[2]{colframe=#1,colback=#1!7,coltitle=white,colbacktitle=#1,fonttitle=\bfseries,title={#2},arc=3pt,boxrule=1pt}
\newtcolorbox{qbox}[1]{colframe=qorange,colback=qbg,coltitle=white,colbacktitle=qorange,fonttitle=\bfseries,title={#1},arc=3pt,breakable,boxrule=1pt}
\newtcolorbox{keybox}{colframe=keygray,colback=keybg,coltitle=white,colbacktitle=keygray,fonttitle=\bfseries,title={$\checkmark$\ Answer Key},arc=3pt,breakable,boxrule=1pt}
\newtcolorbox{recapbox}{colframe=keygray,colback=keybg,coltitle=white,colbacktitle=keygray,fonttitle=\bfseries,title={Question Recap},arc=3pt,breakable,boxrule=1pt}
\newtcolorbox{ideabox}{colframe=titleblue,colback=titleblue!6,coltitle=white,colbacktitle=titleblue,fonttitle=\bfseries,title={The Big Ideas},arc=3pt,breakable,boxrule=1.2pt}
\newtcolorbox{learnbox}{colframe=iagreen,colback=iagreen!5,coltitle=white,colbacktitle=iagreen,fonttitle=\bfseries,title={Concept Essentials},arc=3pt,breakable,boxrule=1.2pt}
\newcommand{\lh}[1]{\par\smallskip{\bfseries\color{iagreen}#1}\par\nobreak\smallskip}

\newcommand{\soln}{\par\smallskip\textit{\textbf{Solution.}}\ }
\newcommand{\workspace}[1]{\par\vspace{3pt}\fcolorbox{black!30}{white}{\begin{minipage}[t][#1][t]{\dimexpr\linewidth-2\fboxsep\relax}\ \end{minipage}}\par}
\newcommand{\sech}[1]{\par\vspace{8pt}{\Large\bfseries\color{iagreen}#1}\par\nobreak\vspace{2pt}{\color{black!20}\hrule height1pt}\vspace{6pt}}

% A compact coordinate plot sized for an example box: \eplot{xmin}{xmax}{ymin}{ymax}{body}
\newcommand{\eplot}[5]{\begin{center}\begin{tikzpicture}
\begin{axis}[width=6.8cm,height=4.4cm,axis lines=middle,xlabel={\tiny$x$},ylabel={\tiny$y$},
grid=both,grid style={gray!16},xmin=#1,xmax=#2,ymin=#3,ymax=#4,
xtick distance=2,ytick distance=2,tick label style={font=\tiny},axis line style={-},samples=120]
#5
\end{axis}\end{tikzpicture}\end{center}}

% A sine-curve plot (degrees on x, 0..360): \splot{ymin}{ymax}{body}
\newcommand{\splot}[3]{\begin{center}\begin{tikzpicture}
\begin{axis}[width=8.6cm,height=4.6cm,axis lines=middle,xlabel={\tiny$x\,(^\circ)$},ylabel={\tiny$y$},
grid=both,grid style={gray!16},xmin=0,xmax=360,ymin=#1,ymax=#2,
xtick distance=90,ytick distance=1,tick label style={font=\tiny},axis line style={-},samples=180]
#3
\end{axis}\end{tikzpicture}\end{center}}

% A small coordinate plot: \plot{xmin}{xmax}{ymin}{ymax}{body}
\newcommand{\plot}[5]{\begin{center}\begin{tikzpicture}
\begin{axis}[width=8.4cm,height=6cm,axis lines=middle,xlabel={\scriptsize$x$},ylabel={\scriptsize$y$},
grid=both,grid style={gray!16},xmin=#1,xmax=#2,ymin=#3,ymax=#4,
xtick distance=2,ytick distance=2,tick label style={font=\tiny},axis line style={-}]
#5
\end{axis}\end{tikzpicture}\end{center}}

% Right triangle (angle theta at bottom-left, right angle at bottom-right):
%   \rtri{drawBase}{drawHeight}{oppLabel}{adjLabel}{hypLabel}
\newcommand{\rtri}[5]{\begin{center}\begin{tikzpicture}[scale=0.85]
\coordinate (A) at (0,0);\coordinate (B) at (#1,0);\coordinate (C) at (#1,#2);
\draw[thick,exblue] (A)--(B)--(C)--cycle;
\draw[black] ($(B)+(-0.32,0)$)--($(B)+(-0.32,0.32)$)--($(B)+(0,0.32)$);
\node[below] at ($(A)!0.5!(B)$){#4};
\node[right] at ($(B)!0.5!(C)$){#3};
\node[above left] at ($(A)!0.5!(C)$){#5};
\node at (0.7,0.22){$\theta$};
\end{tikzpicture}\end{center}}

% General (oblique) triangle with vertex labels and opposite-side labels:
%   \gtri{Alabel}{Blabel}{Clabel}{aSide}{bSide}{cSide}
\newcommand{\gtri}[6]{\begin{center}\begin{tikzpicture}[scale=0.85]
\coordinate (A) at (0,0);\coordinate (B) at (5,0);\coordinate (C) at (1.6,3);
\draw[thick,exblue] (A)--(B)--(C)--cycle;
\node[below left] at (A){#1};\node[below right] at (B){#2};\node[above] at (C){#3};
\node[below] at ($(A)!0.5!(B)$){#6};
\node[right] at ($(B)!0.5!(C)$){#4};
\node[left] at ($(A)!0.5!(C)$){#5};
\end{tikzpicture}\end{center}}

\fancyhead[R]{\footnotesize MCR3U \textperiodcentered{} 6.5}
\begin{document}

\begin{center}
{\LARGE\bfseries\color{titleblue}Worksheet 6.5 --- Pascal's Triangle \& the Binomial Theorem}\\[2pt]
{\small Grade 11 Functions, University (MCR3U) \textperiodcentered{} Unit 6: Sequences and Series}
\end{center}

\begin{ideabox}
Pascal's triangle gives the coefficients for expanding $(a+b)^n$.
\begin{itemize}[leftmargin=*,itemsep=2pt]
  \item Each entry is the sum of the two above it.
  \item Row $n$ gives the coefficients of $(a+b)^n$.
  \item Power of $a$ falls; power of $b$ rises.
\end{itemize}
\end{ideabox}

\begin{learnbox}
\lh{The triangle}Each entry of \textbf{Pascal's triangle} is the sum of the two above it, and row $n$ holds the coefficients of $(a+b)^n$.
\lh{The binomial theorem}$(a+b)^n=\displaystyle\sum_{r=0}^{n}\dbinom{n}{r}a^{n-r}b^{r}$, where $\dbinom{n}{r}$ is a triangle entry and the exponents sum to $n$.
\lh{Single terms}The general term $\dbinom{n}{r}a^{n-r}b^{r}$ lets you find one coefficient — say of $x^3$ — without the full expansion.
\end{learnbox}

\sech{Worked Examples}

\begin{example}{ex0}{Example 1 --- A row}
Write row $4$.\soln Add adjacent entries of row 3 ($1,3,3,1$): $1,\ 1{+}3,\ 3{+}3,\ 3{+}1,\ 1$. \\ Row 4 is $1,4,6,4,1$.
\end{example}

\begin{example}{ex1}{Example 2 --- Cube}
Expand $(a+b)^3$.\soln Coefficients $1,3,3,1$ with falling/rising powers: $a^3+3a^2b+3ab^2+b^3$.
\end{example}

\begin{example}{ex2}{Example 3 --- With numbers}
Expand $(x+2)^3$.\soln Use $1,3,3,1$ with $b=2$: $x^3+3x^2(2)+3x(4)+8=x^3+6x^2+12x+8$.
\end{example}

\begin{example}{ex3}{Example 4 --- Fourth power}
Expand $(x+1)^4$.\soln Row 4 ($1,4,6,4,1$): $x^4+4x^3+6x^2+4x+1$.
\end{example}

\begin{example}{ex4}{Example 5 --- A coefficient}
Coefficient of $a^2b^2$ in $(a+b)^4$?\soln The middle entry of row 4: $6$.
\end{example}

\begin{example}{ex5}{Example 6 --- A row}
Write row $3$.\soln From row 2 ($1,2,1$): $1,3,3,1$.
\end{example}

\begin{example}{ex6}{Example 7 --- Square}
Expand $(a+b)^2$.\soln Coefficients $1,2,1$: $a^2+2ab+b^2$.
\end{example}

\begin{example}{ex7}{Example 8 --- With numbers}
Expand $(x+1)^3$.\soln $1,3,3,1$ with $b=1$: $x^3+3x^2+3x+1$.
\end{example}

\begin{example}{ex8}{Example 9 --- Square}
Expand $(x+3)^2$.\soln $x^2+2(3)x+9=x^2+6x+9$.
\end{example}

\clearpage
\sech{Your Turn --- Show Your Work}
For each question, show all steps clearly in the space provided.

\begin{qbox}{Question 1}
Write row $3$ of Pascal's triangle.\workspace{2.4cm}
\end{qbox}

\begin{qbox}{Question 2}
Expand $(a+b)^2$.\workspace{2.4cm}
\end{qbox}

\begin{qbox}{Question 3}
Expand $(x+1)^3$.\workspace{2.4cm}
\end{qbox}

\begin{qbox}{Question 4}
Expand $(x+3)^2$.\workspace{2.4cm}
\end{qbox}

\begin{qbox}{Question 5}
Coefficient of $x^3$ in $(x+1)^4$?\workspace{2.4cm}
\end{qbox}

\begin{qbox}{Question 6}
Write row $5$ of Pascal's triangle.\workspace{2.4cm}
\end{qbox}

\begin{qbox}{Question 7}
Expand $(a+b)^4$.\workspace{2.4cm}
\end{qbox}

\begin{qbox}{Question 8}
Expand $(x+2)^3$.\workspace{2.4cm}
\end{qbox}

\begin{qbox}{Question 9}
Expand $(x-1)^3$.\workspace{2.4cm}
\end{qbox}

\begin{qbox}{Question 10}
Coefficient of $x^2$ in $(x+1)^4$?\workspace{2.4cm}
\end{qbox}

\begin{qbox}{Question 11}
What do the entries in row $4$ sum to?\workspace{2.4cm}
\end{qbox}

\begin{qbox}{Question 12}
Expand $(x+1)^2$.\workspace{2.4cm}
\end{qbox}

\begin{qbox}{Question 13 --- Challenge}
Expand $(2x+1)^3$.\workspace{3cm}
\end{qbox}

\clearpage
\sech{Question Recap \& Answer Key}

\begin{recapbox}
\begin{multicols}{2}
\small
\begin{enumerate}[leftmargin=*,itemsep=3pt]
  \item Write row $3$ of Pascal's triangle.
  \item Expand $(a+b)^2$.
  \item Expand $(x+1)^3$.
  \item Expand $(x+3)^2$.
  \item Coefficient of $x^3$ in $(x+1)^4$?
  \item Write row $5$ of Pascal's triangle.
  \item Expand $(a+b)^4$.
  \item Expand $(x+2)^3$.
  \item Expand $(x-1)^3$.
  \item Coefficient of $x^2$ in $(x+1)^4$?
  \item What do the entries in row $4$ sum to?
  \item Expand $(x+1)^2$.
  \item Expand $(2x+1)^3$.
\end{enumerate}
\end{multicols}
\end{recapbox}

\begin{keybox}
\begin{multicols}{2}
\small
\begin{enumerate}[leftmargin=*,itemsep=3pt]
  \item $1,3,3,1$
  \item $a^2+2ab+b^2$
  \item $x^3+3x^2+3x+1$
  \item $x^2+6x+9$
  \item $4$
  \item $1,5,10,10,5,1$
  \item $a^4+4a^3b+6a^2b^2+4ab^3+b^4$
  \item $x^3+6x^2+12x+8$
  \item $x^3-3x^2+3x-1$
  \item $6$
  \item $16$
  \item $x^2+2x+1$
  \item $8x^3+12x^2+6x+1$
\end{enumerate}
\end{multicols}
\end{keybox}

\end{document}
